% !PRINTED 411
\weberpage{411}
\begin{center}
§. 132.

\gesperrt{Primitive Roots of Unity.}
\end{center}

That primitive roots of unity exist for every degree $n$ we have not yet proved in the
preceding section. We must first make good this omission, and in doing so we shall also
determine the exact number of the primitive roots of unity.

Let the degree $n$ be decomposed into two factors $a, b$ which are relatively prime to one
another, so that

\[
n = ab,
\]
and let $\alpha$ be an $a$th and $\beta$ a $b$th root of unity. Then the product

\weq{1}{r = \alpha\beta}
is an $n$th root of unity. If $\alpha', \beta'$ are two other $a$th and $b$th roots of
unity, then $r' = \alpha'\beta'$ is also an $n$th root of unity, and it is to be shown
that $r'$ is different from $r$ unless at the same time $\alpha = \alpha'$ and
$\beta = \beta'$.

For since $a, b$ are relatively prime, one can by \S.\,118 determine the integers $x, y$
so that

\weq{2}{ax + by = 1}
holds; and then it follows from (1) that

\[
\alpha = r^{by}, \quad \beta = r^{ax}.
\]
Accordingly $\alpha$ and $\beta$ are completely determined by $r$, and if $\alpha'\beta'$
is also to be equal to $r$, then we must have $\alpha = \alpha'$, $\beta = \beta'$.

If, therefore, we let $\alpha$ in (1) run through all the $a$th, and $\beta$ through all
the $b$th roots of unity, then $r$ takes exactly $ab = n$ different values, and it
follows:

I.~\gesperrt{that in the form $\alpha\beta$ all the $n$th roots of unity can be
represented,}

and further:

II.~\gesperrt{that $r$ is a primitive $n$th root of unity if and only if $\alpha$ is a
primitive $a$th and $\beta$ a primitive $b$th root of unity.}

For, first, let $\mu$ be the smallest positive exponent for which $r^\mu = 1$; then also
$\alpha^\mu \beta^\mu = 1$, and from this it follows by (2), if we raise both sides to the
powers $by$ and $ax$,

\[
\alpha^\mu = 1, \quad \beta^\mu = 1.
\]
